%% appC --- Macierz narzedzi
%%
%% Napisany we wrzesniu 2026. Kazda wersja tutaj drukuje sie z makra
%% przypietej wersji w preamble.tex i zadna nie jest wpisana: ta strona
%% jest jedynym miejscem, gdzie czytelnik widzi wszystkie przypiecia naraz,
%% a przypiecie przepisane tutaj byloby druga kopia faktu, ktory preambula
%% juz posiada. tools/check_structure.py --tools oblewa, gdy makro
%% przypiecia istnieje, a ten dodatek go nie drukuje, oraz gdy wersja jest
%% wpisana zamiast wydrukowana z makra.

\chapter{Macierz narzędzi}
\label{app:tools}

\noindent
Kolumna wersji nie jest wpisana. Każda liczba na tej stronie to to samo
makro, które drukuje strona tytułowa i rozdziały, czyli ten sam napis,
który przypina \code{code/pyproject.toml} i rozwiązuje \code{uv.lock}
\dash{} więc macierz nie może się nie zgadzać z tym, co zainstalowało CI, a
podniesienie wersji rusza tę stronę bez niczyjej edycji. Bramka oblewa
build, jeśli przypięcie istnieje, a ta tabela go nie pokazuje.

Czytaj pierwszą kolumnę luźno. Część z nich to dokładne odpowiedniki, a
część tylko najbliższa rzecz, którą inżynier .NET już zna; tam, gdzie oba
różnią się w sposób, który ma znaczenie, rozdział wskazany po prawej jest
tym, w którym ta różnica jest wydana.

\section{Język i jego narzędziówka}
\label{sec:appC-toolchain}

\begin{toolmatrix}{.NET}{Python}{Sprawdzone na}{Rozdz.}
środowisko uruchomieniowe & CPython & \pypatch & \ref{ch:runtime} \\
CLI \code{dotnet}, NuGet & \pkg{uv} & \uvver & \ref{ch:packaging} \\
\code{dotnet format}, analizatory & \pkg{ruff} & \ruffver &
  \ref{ch:packaging} \\
sprawdzanie typów przez kompilator & \pkg{pyright}, ściśle &
  \pyrightver & \ref{ch:typing} \\
--- & \pkg{mypy}, druga opinia & \mypyver & \ref{ch:typing} \\
\end{toolmatrix}

\noindent
Przypięte są dwa checkery, a nie jeden, i rozdział~\ref{ch:typing} mierzy,
dlaczego: na tych samych czterech funkcjach pyright nie zgłasza nic, a mypy
zgłasza dwa, bo jeden stosuje system typów, a drugi dokłada opinię, której
system typów nie ma.

\section{Wdrażanie}
\label{sec:appC-shipping}

\begin{toolmatrix}{.NET}{Python}{Sprawdzone na}{Rozdz.}
wiązanie modelu, walidacja & \pkg{pydantic} & \pydanticver &
  \ref{ch:typing} \\
\code{IOptions<T>} & \pkg{pydantic-settings} & \pysettingsver &
  \ref{ch:services} \\
ASP.NET Core & \pkg{fastapi} & \fastapiver & \ref{ch:services} \\
Kestrel & \pkg{uvicorn} & \uvicornver & \ref{ch:services} \\
\code{HttpClient} & \pkg{httpx} & \httpxver & \ref{ch:asyncio} \\
Entity Framework & \pkg{sqlalchemy} & \sqlalchemyver & \ref{ch:data} \\
migracje EF & \pkg{alembic} & \alembicver & \ref{ch:data} \\
--- & \pkg{aiosqlite}, do listingów asynchronicznych & \aiosqlitever &
  \ref{ch:data} \\
DataTable & \pkg{polars} & \polarsver & \ref{ch:data} \\
\end{toolmatrix}

\section{Testowanie i eksploatacja}
\label{sec:appC-testing}

\begin{toolmatrix}{.NET}{Python}{Sprawdzone na}{Rozdz.}
xUnit & \pkg{pytest} & \pytestver & \ref{ch:testing} \\
wsparcie testów asynchronicznych & \pkg{pytest-asyncio} & \pytestasyncver &
  \ref{ch:testing} \\
FsCheck & \pkg{hypothesis} & \hypothesisver & \ref{ch:testing} \\
coverlet & \pkg{coverage} & \coveragever & \ref{ch:testing} \\
\code{ILogger<T>} & \pkg{structlog} & \structlogver &
  \ref{ch:observability} \\
OpenTelemetry & \pkg{opentelemetry-sdk} & \otelver &
  \ref{ch:observability} \\
\code{dotnet list package -{}-vulnerable} & \pkg{pip-audit} &
  \pipauditver & \ref{ch:observability} \\
\end{toolmatrix}

\section{Zestaw do AI}
\label{sec:appC-ai}

\begin{toolmatrix}{.NET}{Python}{Sprawdzone na}{Rozdz.}
SDK & \pkg{anthropic} & \anthropicver & \ref{ch:aitoolkit} \\
SDK & \pkg{openai} & \openaiver & \ref{ch:aitoolkit} \\
\end{toolmatrix}

\begin{note}
Oba SDK zależą od dystrybucji \pkg{httpx2}, a nie od \pkg{httpx}
przypiętego powyżej. To różne pakiety, które dzielą pochodzenie i zapis, i
dlatego sprawdzenie \api{isinstance} względem tego, który przypiąłeś,
zawodzi na kliencie, który jest oczywiście klientem httpx.
Rozdział~\ref{ch:aitoolkit} to mierzy.
\end{note}

\section{Budowanie książki}
\label{sec:appC-book}

\noindent
Dwa przypięcia nie są w tabelach powyżej, bo nie są bibliotekami, które
instaluje czytelnik. Interpreterem, na który celuje książka, jest
Python~\pyver, a dokładną budową, na której działał każdy listing, jest
\pypatch, czyli to, co nazywa \code{code/.python-version}. A własny C\#
książki \dash{} porównania, które robią jej rozdziały, i rozwiązania w
dodatku~\ref{app:interview} \dash{} jest kompilowany i testowany w CI
względem .NET~\dotnetver, przypiętego raz w
\code{code/csharp/global.json} i stamtąd czytanego przez każdy workflow,
który go potrzebuje.
